\documentclass[10pt,a4paper]{amsart}
\usepackage[margin=19mm]{geometry}
\usepackage{amsmath,amssymb,mathtools}
\usepackage[scr=rsfs,cal=euler]{mathalfa}
\usepackage{xfrac}
\usepackage[unicode,hidelinks,bookmarks=false]{hyperref}
\newcommand{\ie}{i.e.\ }
\newcommand{\cf}{cf.\ }
\newcommand{\resp}{respectively\ }
\newcommand{\Subsection}{\subsection}
\providecommand{\nobreakdash}{\nobreak\hbox{-}}
\setlength{\emergencystretch}{2em}
\multlinegap0pt
\usepackage{amssymb}
\usepackage{bm}
\usepackage{mathtools}
\usepackage{braket}
\newtheorem{lemma}{Lemma}
\newcommand{\R}{\mathbb{R}}
\newcommand{\one}{\bm{1}}
\title{An Effective SPDE Model for a Schr\"odinger Equation with a Fluctuating Magnetic Potential}
\author{Franco Flandoli, Reika Fukuizumi, Francesco Grotto, Eliseo Luongo}
\date{Source: arXiv:2609.03556v1 (CC BY 4.0)}
\begin{document}
\maketitle

\noindent\emph{Source and licence.} An Effective SPDE Model for a Schr\"odinger Equation with a Fluctuating Magnetic Potential. Version arXiv:2609.03556v1 (CC BY 4.0), distributed by arXiv under the Creative Commons Attribution 4.0 International licence, http://creativecommons.org/licenses/by/4.0/, as recorded in the arXiv licence metadata for this exact version and shown on the abstract page https://arxiv.org/abs/2609.03556v1. Full licence notice: \texttt{licenses/arxiv-2609.03556-CC-BY-4.0.txt}.\par\smallskip

\noindent\textit{Editorial scope.} Counted unit: the article's lemma lem:limit\_cov (source0.tex lines 3073--3078). The article's complete original proof of it is delivered with it (source0.tex lines 3079--3127, one \texttt{proof} environment of 49 lines). No mathematics is written, repaired or re-derived: the statement and the proof are the article's own lines, in the article's own notation, and no retained line is substituted or re-flowed.  No further source-local context is needed: every cross-reference of every retained line resolves inside the two delivered spans.  Nothing was omitted from any retained span, so the package carries no omission record.  No retained line contains a citation, so the wrapper carries no bibliography and no bibliography entry.  No retained line contains a figure environment, an image-inclusion command or drawing code, so no image is delivered and the package declares no asset dependency.  Editorial formatting only, in addition: an AMS \texttt{amsart} preamble that restates the article's own declarations and macros used by the retained lines, namely the environments \texttt{lemma} on separate counters, together with the standard packages those lines need.  The retained environments print in this document as Lemma 1 in order of appearance, whereas the article prints the same items with the numbering of its own full document.  The complete native source of the article as distributed in the arXiv e-print for this exact version is bundled unchanged as \texttt{sources/209306/source0.tex}.
\begin{lemma}\label{lem:limit_cov}
    For each $\chi\in C^{\infty}_c((0,T)\times \R^d)$ it holds
    \begin{align*}
        \lim_{\tau\rightarrow 0}\mathbb{E}\left[\left|\int_0^T \left\langle \frac{|A^\tau|^2-V}{\sqrt{\tau}}\bar\psi_t,\chi_t\right\rangle dt \right|^2 \right]=2 \sum_{k,j\in I}\int_0^T |\langle \sigma_k\cdot\sigma_j\bar\psi_t, \chi_t \rangle|^2 dt.
    \end{align*}
\end{lemma}

\begin{proof}
Let us denote by $A^{k,\tau}_t=\frac{\sqrt{2}}{\sqrt{\tau}}\int_{-\infty}^t e^{-\frac{(t-r)}{\tau}}dW^{k}_r$. Therefore
\begin{align*}
    A^{\tau}_t=\sum_{k\in I } A^{k,\tau}_t\sigma_k\bar\psi_t
\end{align*}
and 
\begin{align}
\frac{1}{\tau}&\mathbb{E}\left[\left|\int_0^T \langle |A^\tau_t|^2-V, \chi_t\bar{\psi}_t\rangle dt\right|^2\right]\notag
\\ & = \frac{1}{\tau}\mathbb{E}\left[\int_0^T\int_0^t \langle |A^\tau_t|^2-V, \chi_t \bar{\psi}_t\rangle  
\langle |A^\tau_s|^2-V, \chi_s\bar{\psi}^*_s\rangle ds dt\right]\notag\\ &+\frac{1}{\tau}\mathbb{E}\left[\int_0^T\int_0^t \langle |A^\tau_t|^2-V, \chi_t \bar{\psi}_s\rangle  
\langle |A^\tau_s|^2-V, \chi_s\bar{\psi}^*_t\rangle ds dt\right]\notag\\
&= \frac{1}{\tau}\sum_{k,j\in I}\int_0^T\int_0^t 
\mathbb{E}\left[(|A^{k,\tau}_t|^2-1)(|A^{j,\tau}_s|^2-1)\right]
\langle |\sigma_k|^2, \chi_t\bar{\psi}_t\rangle 
\langle |\sigma_j|^2, \chi_s\bar{\psi}^*_s\rangle ds dt\notag\\&
+\frac{1}{\tau}\sum_{k,j\in I}\int_0^T\int_0^t 
\mathbb{E}\left[(|A^{k,\tau}_t|^2-1)(|A^{j,\tau}_s|^2-1)\right]
\langle |\sigma_k|^2, \chi_t\bar{\psi}_s\rangle 
\langle |\sigma_j|^2, \chi_s\bar{\psi}^*_t\rangle ds dt \notag\\
& +\frac{1}{\tau}\sum_{\substack{k,j,k',j'\in I\\ k'\neq k\\ j'\neq j}}
\int_0^T\int_0^t 
\mathbb{E}\left[A^{k,\tau}_tA^{k',\tau}_tA^{j,\tau}_s A^{j',\tau}_s \right] 
\langle \sigma_k \cdot\sigma_{k'}, \chi_t\bar{\psi}_t\rangle 
\langle \sigma_j\cdot\sigma_{j'}, \chi_s\bar{\psi}^*_s\rangle ds dt \notag\\
& + \frac{1}{\tau}\sum_{\substack{k,j,k',j'\in I\\ k'\neq k\\ j'\neq j}}
\int_0^T\int_0^t 
\mathbb{E}\left[A^{k,\tau}_tA^{k',\tau}_tA^{j,\tau}_s A^{j',\tau}_s \right] 
\langle \sigma_k \cdot\sigma_{k'}, \chi_t\bar{\psi}_s\rangle 
\langle \sigma_j\cdot\sigma_{j'}, \chi_s\bar{\psi}^*_t\rangle ds dt\notag\\ 
&=: I^{1,\tau}+I^{2,\tau}.
\label{eq:rewriting}
\end{align}
It also holds 
\begin{align}
  \mathbb{E}\left[(|A^{k,\tau}_t|^2-1)(|A^{j,\tau}_s|^2-1)\right]&=2\delta_{k,j} e^{-\frac{2(t-s)}{\tau}}\label{eq_1_expt} \\ 
\one_{k\neq k', j\neq j'}\mathbb{E}\left[A^{k,\tau}_tA^{k',\tau}_tA^{j,\tau}_s A^{j',\tau}_s \right]&=\one_{k\neq k', j\neq j'}\left(\delta_{k,j}\delta_{k',j'}+\delta_{k,j'}\delta_{k',j}\right)\label{eq_2_expt}e^{-\frac{2(t-s)}{\tau}}.
\end{align}
Combining \eqref{eq:rewriting}, \eqref{eq_1_expt} and \eqref{eq_2_expt} we get
\begin{align*}
 \frac{1}{\tau}\mathbb{E}\left[\left(\int_0^T \langle |A^\tau_t|^2-V, \chi_t\bar{\psi}_t\rangle dt\right)^2\right]&=\frac{2}{\tau}\sum_{k\in I}\int_0^T\int_0^t 
e^{-\frac{2(t-s)}{\tau}}
\langle |\sigma_k|^2, \chi_t\bar{\psi}_t\rangle 
\langle |\sigma_k|^2, \chi_s \bar{\psi}^*_s\rangle ds dt\\ & +\frac{2}{\tau}\sum_{k\in I}\int_0^T\int_0^t 
e^{-\frac{2(t-s)}{\tau}}
\langle |\sigma_k|^2, \chi_t\bar{\psi}_s\rangle 
\langle |\sigma_k|^2, \chi_s \bar{\psi}^*_t\rangle ds dt\\ & +\frac{2}{\tau}\sum_{\substack{k,j\in I\\ k\neq j}}\int_0^T \int_0^t e^{-\frac{2(t-s)}{\tau}}\langle \sigma_k \cdot\sigma_j,\chi_t\bar{\psi}_t\rangle \langle \sigma_k\cdot \sigma_j,\chi_s\bar{\psi}^*_s\rangle ds dt\\ &+\frac{2}{\tau}\sum_{\substack{k,j\in I\\ k\neq j}}\int_0^T \int_0^t e^{-\frac{2(t-s)}{\tau}}\langle \sigma_k \cdot\sigma_j,\chi_t\bar{\psi}_s\rangle \langle \sigma_k\cdot \sigma_j,\chi_s\bar{\psi}_t^*\rangle ds dt\\ & \rightarrow 2 \sum_{k,j\in I}\int_0^T |\langle \sigma_k\cdot\sigma_j \bar{\psi}_t, \chi_t \rangle|^2dt\quad\text{as }\tau\rightarrow 0
\end{align*}
which is the claim.
\end{proof}

\end{document}
